% !PRINTED 227
\weberpage{227}
The expression on the right-hand side of (7) is a \gesperrt{bilinear}
function of the $t$ and $\tau$; when developed, let it have the form
\[
\sum_{0,\,n-2}^{h} \sum_{0,\,n-2}^{k} B_{h,k} t_h \tau_k,
\]
with the condition $B_{i,k} = B_{k,i}$. The \gesperrt{Bezoutiant} is then
obtained by setting $\tau = t$, thus
\weq{8}{B = \sum_{0,\,n-2}^{h,k} B_{h,k} t_h t_k,}
in which the coefficients $B_{i,k}$ are quadratic functions of the $a$;
and from (7) it follows
\weq{9}{S(y z) - \frac{1}{n} S(y) S(z) = \sum_{0,\,n-2}^{h,k} B_{h,k} t_h \tau_k,}
\weq{10}{S(y^2) = \frac{1}{n} [S(y)]^2 + B.}
By formula (8), \S.\,68, we have
\[
\begin{aligned}
y - \frac{1}{n} S(y) &= t_{n-2} F_0 + t_{n-3} F_1 + \cdots + t_0 F_{n-2} \\
z - \frac{1}{n} S(z) &= \tau_{n-2} F_0 + \tau_{n-3} F_1 + \cdots + \tau_0 F_{n-2},
\end{aligned}
\]
and if we multiply the two together,
\[
y z - \frac{1}{n} [y S(z) + z S(y)] + \frac{1}{n^2} S(y) S(z)
= \sum_{0,\,n-2}^{h,k} \tau_h t_k F_{n-h-2} F_{n-k-2}.
\]
If we sum this formula over all the roots $x$, that is, take
the sum denoted by $S$, it follows that
\[
S(y z) - \frac{1}{n} S(y) S(z) = \sum^{h,k} \tau_h t_k S(F_{n-h-2} F_{n-k-2}),
\]
and comparison with (9) shows
\weq{11}{B_{h,k} = S(F_{n-h-2} F_{n-k-2})}
We denote by $\varDelta$ the determinant of the function $B$,
and thus set
\[
\varDelta = \sum \pm B_{0,0} B_{1,1} \ldots B_{n-2,\,n-2}.
\]
By the multiplication theorem for determinants, and having
regard to the formulas
\[
S F_0 = 0, \quad S F_1 = 0, \ldots S F_{n-2} = 0
\]
%CONT>>
